\documentclass[aspectratio=169]{beamer}

\usetheme{colorful-dream}

\usepackage{minted}
\setminted{fontsize=\footnotesize, breaklines=true}

\title{Basics of Python for Data Science}
\subtitle{Lecture 2: NumPy, Matplotlib, and Engineering Signals}
\author{Lu Pa}
\date{\today}

% --- Title Slide Logos ---
\titlegraphic{
  \vspace{0.5cm}
  \begin{center}
    \includegraphics[height=1.8cm]{logo_university_placeholder.png}
    \hspace{1cm}
    \includegraphics[height=1.8cm]{logo_faculty_placeholder.png}
  \end{center}
}

\begin{document}

% ---------------------------------------------------------
\begin{frame}
  \titlepage
\end{frame}

% ---------------------------------------------------------
\section{PART I}
\begin{frame}[plain]
  \begin{center}
    \vspace{2cm}
    {\usebeamerfont{title}\color{white}\Huge PART I}
  \end{center}
\end{frame}

% ---------------------------------------------------------
\begin{frame}{From Python to Scientific Computing}
  \begin{itemize}
    \item Python becomes powerful when combined with scientific libraries.
    \item Today we focus on:
      \begin{itemize}
        \item \textbf{NumPy} — numerical arrays and vectorization
        \item \textbf{Matplotlib} — engineering plots
      \end{itemize}
    \item These two libraries form the foundation of scientific Python.
  \end{itemize}
\end{frame}

% ---------------------------------------------------------
\begin{frame}{What is NumPy?}
  \begin{itemize}
    \item Core library for numerical computing in Python.
    \item Provides:
      \begin{itemize}
        \item fast multidimensional arrays
        \item vectorized operations
        \item mathematical functions
        \item random number generation
      \end{itemize}
    \item Written in C for performance.
    \item Backbone of Pandas, SciPy, scikit-learn, and more.
  \end{itemize}
\end{frame}

% ---------------------------------------------------------
\begin{frame}[fragile]{Creating NumPy Arrays}
\begin{minted}{python}
import numpy as np

a = np.array([1, 2, 3])
b = np.linspace(0, 10, 100)
c = np.zeros(50)
\end{minted}
\end{frame}

% ---------------------------------------------------------
\begin{frame}[fragile]{Vectorized Operations}
\begin{minted}{python}
v = np.array([10, 20, 30])
v_ms = v / 3.6
\end{minted}

\begin{itemize}
  \item Operations apply to entire arrays at once.
  \item Much faster and cleaner than Python loops.
\end{itemize}
\end{frame}

% ---------------------------------------------------------
\begin{frame}{Engineering Example: Speed Signal}
  \begin{itemize}
    \item Extract arrays from the DataFrame:
      \begin{itemize}
        \item time
        \item speed
      \end{itemize}
    \item Compute:
      \begin{itemize}
        \item acceleration
        \item jerk
        \item smoothed speed
      \end{itemize}
  \end{itemize}
\end{frame}

% ---------------------------------------------------------
\begin{frame}[fragile]{Extracting Arrays from DataFrame}
\begin{minted}{python}
t = df["time_s"].values
v = df["speed_kmh"].values
\end{minted}
\end{frame}

% ---------------------------------------------------------
\begin{frame}[fragile]{Acceleration via Finite Differences}

\begin{equation}
a(t) \approx \frac{\Delta v}{\Delta t}
\end{equation}

\begin{minted}{python}
dt = np.diff(t)
dv = np.diff(v)
accel = dv / dt
\end{minted}
\end{frame}

% ---------------------------------------------------------
\begin{frame}[fragile]{Jerk (Derivative of Acceleration)}

\begin{equation}
j(t) \approx \frac{\Delta a}{\Delta t}
\end{equation}

\begin{minted}{python}
jerk = np.diff(accel) / dt[:-1]
\end{minted}
\end{frame}

% ---------------------------------------------------------
\begin{frame}[fragile]{Smoothing with Moving Average}
\begin{minted}{python}
window = 15
v_smooth = np.convolve(v, np.ones(window)/window, mode="same")
\end{minted}
\end{frame}

% ---------------------------------------------------------
\begin{frame}{Matplotlib for Engineering Plots}
  \begin{itemize}
    \item Flexible and powerful plotting library.
    \item Ideal for:
      \begin{itemize}
        \item time series
        \item multi-panel dashboards
        \item histograms
        \item scatter plots
      \end{itemize}
    \item Highly customizable.
  \end{itemize}
\end{frame}

% ---------------------------------------------------------
\section{PART II}
\begin{frame}[plain]
  \begin{center}
    \vspace{2cm}
    {\usebeamerfont{title}\color{white}\Huge PART II}
  \end{center}
\end{frame}

% ---------------------------------------------------------
\begin{frame}[fragile]{Multi-panel Engineering Plot}
\begin{minted}{python}
fig, ax = plt.subplots(2, 2, figsize=(12, 8))

ax[0,0].plot(t, v)
ax[0,0].set_title("Speed [km/h]")

ax[0,1].plot(t[:-1], accel)
ax[0,1].set_title("Acceleration [m/s²]")

ax[1,0].plot(t, v_smooth)
ax[1,0].set_title("Smoothed Speed")

ax[1,1].plot(t, df["engine_rpm"])
ax[1,1].set_title("Engine RPM")

for a in ax.ravel():
    a.grid(True)

plt.tight_layout()
plt.show()
\end{minted}
\end{frame}

% ---------------------------------------------------------
\begin{frame}{Practice Tasks}
  \begin{itemize}
    \item Compute jerk and visualize it.
    \item Compare finite difference vs \texttt{np.gradient}.
    \item Try different smoothing window sizes.
    \item Plot a histogram of acceleration.
    \item Build a 2×2 dashboard of key signals.
  \end{itemize}
\end{frame}

% ---------------------------------------------------------
\end{document}
